\documentclass[11pt]{article}
\usepackage{amsmath,amssymb,amsthm}
\usepackage[utf8]{inputenc}
\newcommand{\dd}{\mathrm{d}}
\newcommand{\ii}{\mathrm{i}}
\newcommand{\ee}{\mathrm{e}}
\newcommand{\Order}{\mathcal{O}}

\newtheorem{proposition}{Proposition}

\title{Task S.1: Truncation Validity --- When Can $K_4$ Be Neglected?}
\author{HaibiMathAgent}
\date{March 2026}

\begin{document}
\maketitle

\section{Problem Statement}

GPU-Claude-Opus (Task J.1) found that a four-body coupling $K_4$ produces
$\sim12\%$ variation in $r_1$ for $K_4\in[-1,1]$ at $N=50$.  We need to
establish theoretically when truncation at three-body is justified.

\section{OA Framework for Four-Body Terms}

The most general four-body coupling term is:
\begin{equation}
  T_i^{(4)} = \frac{K_4}{N^3}\sum_{j,k,l}\sin(3\theta_j - \theta_k - \theta_l - \theta_i).
\end{equation}

In the mean-field limit, this becomes:
\begin{equation}
  T_i^{(4)} \xrightarrow{N\to\infty} K_4\,\mathrm{Im}\bigl[Z_3\bar{Z}_1^2\,\ee^{-\ii\theta_i}\bigr].
\end{equation}

On the OA manifold, $Z_3 = z^3$, $Z_1 = z$, so:
\begin{equation}
  Z_3\bar{Z}_1^2 = z^3\bar{z}^2 = |z|^4 z = r^4 z.
\end{equation}

The modified radial dynamics become:
\begin{equation}
  \dot{r} = -\Delta r + \frac{r(1-r^2)}{2}\bigl(K_2 + K_3 r^2 + K_4 r^4\bigr).
  \label{eq:r-K4}
\end{equation}

\section{Scaling Argument for Time-Delay Expansion}

From the time-delay equivalence (Task A.1):
\begin{align}
  K_2^{\mathrm{eff}} &= K_2 + \Order(\tau), \\
  K_3^{\mathrm{eff}} &= \frac{K_2^2\tau}{4} + \Order(\tau^2), \\
  K_4^{\mathrm{eff}} &= c_4 K_2^3\tau^2 + \Order(\tau^3),
\end{align}
where $c_4$ is a numerical constant from the $\Order(\tau^2)$ expansion.

\begin{proposition}[Truncation condition]
  The truncation $K_4=0$ is justified when:
  \begin{equation}
    \frac{K_4 r^4}{K_3 r^2 + K_2} \ll 1,
    \label{eq:truncation}
  \end{equation}
  i.e., the four-body contribution to the effective coupling is small
  compared to the two-body and three-body terms.

  For the time-delay case:
  \begin{equation}
    \frac{K_4^{\mathrm{eff}}}{K_3^{\mathrm{eff}}} \sim K_2\tau,
  \end{equation}
  so truncation is valid when $K_2\tau \ll 1$ (small delay-coupling product).
\end{proposition}

\section{Quantitative Estimate}

From the GPU data at $K_2=2$, $K_3=0.5$, $N=50$:
\begin{itemize}
  \item $r_1(K_4=0) \approx 0.888$
  \item $r_1(K_4=\pm 1)$: range $[0.833, 0.940]$
  \item Relative variation: $12\%$
\end{itemize}

At the measured $r^* \approx 0.888$:
\begin{equation}
  \frac{K_4 r^4}{K_3 r^2 + K_2}
  = \frac{1.0 \times 0.888^4}{0.5 \times 0.888^2 + 2.0}
  = \frac{0.622}{2.394} = 0.26.
\end{equation}

This is $\sim 26\%$ --- not negligible, consistent with the observed $12\%$
effect on $r$.  The truncation criterion \eqref{eq:truncation} correctly
predicts that $K_4$ matters at these parameter values.

\section{Regime of Validity}

\begin{itemize}
  \item For $K_3 \lesssim K_2/2$ and $r^* < 0.7$: truncation error $<5\%$.
  \item For $K_3 \sim K_2$ and $r^* > 0.8$: truncation error $>10\%$,
    four-body terms should be included.
  \item In the time-delay framework with $K_2\tau < 0.3$: truncation safe.
\end{itemize}

\section{Impact on Paper}

The three-body truncation is valid for the parameter regime studied in most
of our paper ($K_3 < K_2$, moderate $r^*$).  Near the explosive boundary
($K_3 \approx K_2$, $r^*$ near 1), four-body terms become significant and
our predictions should be treated as qualitative.  This should be stated
explicitly in the paper's Discussion section.

\end{document}
